%% ch10-compliance.tex — Compliance Reference
%% JA4proxy Technical Reference Manual

\chapter{Compliance Reference}
\index{compliance}
\label{ch:compliance}

\begin{chapterquote}
\ja\ processes IP addresses and TLS fingerprints in the course of its security
function. This chapter documents what personal data is processed, for how long,
what controls apply, and how to satisfy data subject requests.
\end{chapterquote}

\newthought{Compliance obligation depends on deployment context.} \ja\ is a
security proxy; the operator who deploys it is the data controller. This chapter
provides the technical facts that a data controller needs to assess their compliance
obligations. It does not constitute legal advice.

\section{Data Inventory}
\index{data inventory}
\index{personal data}

\subsection{What Data is Processed}

\begin{fullwidth}
\begin{table}[h]
\caption{Data elements processed by \ja}
\label{tab:data-inventory}
\begin{tabular}{@{}lllp{4.5cm}@{}}
\toprule
\textbf{Data element} & \textbf{Personal data?} & \textbf{Where stored} & \textbf{Purpose} \\
\midrule
Client IP address (IPv4/IPv6) & Yes (see note) & Redis, logs, Prometheus labels & Pipeline decisions; rate limiting; ban enforcement; CIDR density \\
JA4 fingerprint & No (device, not person) & Redis, logs & TLS client identification \\
JA4X fingerprint & No & Redis, logs & Extended TLS fingerprint \\
JA4T fingerprint & No & Redis, logs & TCP stack fingerprint \\
SNI hostname & No (operator's domain) & Logs & Pipeline signal; not personal data of client \\
Country code & No (GeoIP; not precise location) & Redis, logs & Country-level access control \\
ASN organisation name & No & Logs (transient) & Signal computation \\
AbuseIPDB score & Derivative of IP (see note) & Redis (cache) & Signal computation \\
RDAP organisation & No & Redis (cache) & Signal computation \\
PTR/FCrDNS hostname & Potentially (if ISP assigns personal hostnames) & In-process cache (transient) & Signal computation \\
Beaconing timestamps & Linked to IP (see note) & Redis Sorted Set & Beaconing detection \\
Connection metadata (port, timing) & Potentially linked to IP & Logs, Prometheus & Operational metrics \\
\bottomrule
\end{tabular}
\end{table}
\end{fullwidth}

\begin{notebox}[IP Addresses as Personal Data]
Under GDPR Recital 30, IP addresses are personal data when they can be linked to a
natural person. Dynamic IP addresses assigned by ISPs to residential customers are
generally considered personal data in the EU. Static IP addresses assigned to
organisations are generally not personal data. \ja\ processes both categories and
cannot distinguish them at the time of processing. Operators must account for this
in their DPIA.
\end{notebox}

\subsection{Legal Basis for Processing}

The operator must establish a legal basis for processing IP addresses. Typical bases:

\begin{description}
  \item[Legitimate interests (Art. 6(1)(f))] Security monitoring and DDoS/bot
    mitigation is a legitimate interest of the controller. The processing is
    proportionate because: only metadata is processed (no content decryption),
    data is retained for the minimum period necessary, and the security purpose
    is genuine.
  \item[Legal obligation (Art. 6(1)(c))] If the operator is subject to regulations
    requiring security controls (PCI-DSS, NIS2), processing for compliance purposes
    has a legal basis.
\end{description}

\section{Data Retention}
\index{data retention}
\index{GDPR!data retention}

\subsection{Redis Key TTLs}

\begin{fullwidth}
\begin{table}[h]
\caption{Redis key TTLs for compliance purposes}
\label{tab:ttl-compliance}
\begin{tabular}{@{}llll@{}}
\toprule
\textbf{Key type} & \textbf{Contains IP?} & \textbf{TTL} & \textbf{Notes} \\
\midrule
\code{ban:\{ip\}} & Yes & 300s–604800s & Escalating ban. Max 7 days. \\
\code{rate:\{strategy\}:\{key\}} & Yes (in key) & 60s window & Sliding window; entries expire with window \\
\code{beacon:\{ip\}:\{ja4\}} & Yes (in key) & 30min window & Timestamps expire; key persists until last entry expires \\
\code{conn:\{ip\}} & Yes (in key) & Connection lifetime & Deleted when connection count reaches 0 \\
\code{visit:\{ip\}} & Yes (in key) & None (persistent) & Manual deletion required for erasure \\
\code{sess:\{ip\}} & Yes (in key) & None (persistent) & Manual deletion required for erasure \\
\code{hll:cidr24:\{cidr\}} & Yes (CIDR containing IP) & 30 days & Approximate; cannot extract individual IPs \\
\code{abuseipdb:\{ip\}} & Yes (in key) & 3600s & Cache; auto-expires \\
\code{geo:\{ip\}} & Yes (in key) & 86400s & Cache; auto-expires \\
\code{rdap:\{ip\}} & Yes (in key) & 86400s & Cache; auto-expires \\
\code{findings:\{ip\}} & Yes (in key) & 300s & Analytics output; auto-expires \\
\code{management:policy\_audit} & No & None (max 1000 entries) & Audit log; no IP data \\
\bottomrule
\end{tabular}
\end{table}
\end{fullwidth}

\subsection{Log Retention}

Connection logs are emitted to stdout and shipped to Loki (or another log aggregator).
The retention period is configured in the log aggregation system:

\begin{itemize}
  \item \textbf{Recommended retention:} 30 days for operational logs
  \item \textbf{Security incident retention:} 90 days (check applicable law)
  \item \textbf{Legal hold:} Operator-specific; coordinate with legal team
\end{itemize}

Log retention is configured in Loki's \code{storage\_config} section. \ja\ does not
enforce log retention directly — this is the responsibility of the log aggregation
pipeline.

\subsection{Prometheus Metrics Retention}

Prometheus metrics contain IP addresses only in specific label combinations (notably
the ban metrics, which label by IP). Default Prometheus retention is 15 days.
For compliance, either:
\begin{enumerate}
  \item Reduce retention to 7 days for metrics containing IP labels
  \item Configure remote write to a long-term storage system with access controls
  \item Consider using IP prefix labels (/24 or /48) instead of full IPs for metrics
\end{enumerate}

\section{Right to Erasure (GDPR Art. 17)}
\index{GDPR!right to erasure}
\index{DSAR}

\subsection{Data Subject Access Request (DSAR) Process}

When a data subject submits a right-to-erasure request identifying their IP address,
the operator must delete all personal data associated with that IP. The following
Redis keys must be deleted:

\begin{lstlisting}[language=bash]
# DSAR erasure script for a specific IP address
# Replace TARGET_IP with the subject's IP address

TARGET_IP="203.0.113.42"

# Delete all Redis keys containing this IP
redis-cli DEL "ban:${TARGET_IP}"
redis-cli DEL "conn:${TARGET_IP}"
redis-cli DEL "visit:${TARGET_IP}"
redis-cli DEL "sess:${TARGET_IP}"
redis-cli DEL "abuseipdb:${TARGET_IP}"
redis-cli DEL "geo:${TARGET_IP}"
redis-cli DEL "rdap:${TARGET_IP}"
redis-cli DEL "findings:${TARGET_IP}"

# Delete beaconing keys (need to enumerate JA4 variants)
redis-cli SCAN 0 MATCH "beacon:${TARGET_IP}:*" COUNT 100
# Delete each matched key

# Delete rate limit keys
redis-cli SCAN 0 MATCH "rate:*:${TARGET_IP}*" COUNT 100
# Delete each matched key

# Logs: use log aggregation platform API to delete/redact log entries
# HyperLogLog: individual IPs cannot be removed from HLL — document this limitation
\end{lstlisting}

\begin{warningbox}[HyperLogLog Erasure Limitation]
HyperLogLog data structures (\code{hll:cidr24:} and \code{hll:cidr48:}) do not
support removal of individual elements. Once an IP has been added to a HLL, it cannot
be removed without deleting the entire HLL key (which would also remove all other
IPs from that CIDR). Document this technical limitation in your DPIA and assess
whether HLL data constitutes personal data under your jurisdiction (it does not
reveal the actual IP address, only an approximate count).
\end{warningbox}

\subsection{IP Anonymisation Option}
\index{IP anonymisation}

For environments where storing full IP addresses is not appropriate, \ja\ can be
configured to truncate IP addresses in logs and Redis keys:

\begin{itemize}
  \item IPv4: truncate last octet (e.g.\ \code{203.0.113.42} → \code{203.0.113.0})
  \item IPv6: truncate to /48 (e.g.\ \code{2001:db8::1} → \code{2001:db8::})
\end{itemize}

\begin{dangerbox}[IP Anonymisation Degrades Security Function]
Truncating IP addresses prevents ban enforcement (you cannot ban \code{203.0.113.0/24}
without blocking 255 potentially innocent IPs). It also prevents return visitor
tracking, session resumption ratio, and concurrent connection counting. IP anonymisation
severely degrades the security function of \ja\ and should only be used when legally
required, with operator acknowledgement of the security trade-off.
\end{dangerbox}

\section{PCI-DSS Considerations}
\index{PCI-DSS}

\ja\ contributes to the following PCI-DSS v4.0 requirements when deployed in front
of a cardholder data environment:

\begin{fullwidth}
\begin{table}[h]
\caption{PCI-DSS v4.0 control contributions}
\label{tab:pci-dss}
\begin{tabular}{@{}llp{7cm}@{}}
\toprule
\textbf{Requirement} & \textbf{Control} & \textbf{\ja\ contribution} \\
\midrule
1.2 & Network access controls & Enforces perimeter access control via IP, country, and TLS fingerprint \\
1.3 & Network access restriction & Blocks traffic from known malicious sources (Spamhaus, blacklisted countries) \\
6.4 & WAF/bot protection & TLS fingerprint analysis detects automated tools and scanners \\
6.5 & Change management & Config audit log (\code{management:policy\_audit}) provides change trail \\
10.2 & Audit logs & Structured JSON logs cover all access events with action, IP, and reason \\
10.3 & Audit log protection & Logs shipped to Loki; proxy cannot delete its own logs \\
10.5 & Log retention & Log retention policy (see Section above) \\
12.10 & Incident response | Security monitoring via anomaly detection and alerting \\
\bottomrule
\end{tabular}
\end{table}
\end{fullwidth}

\section{SOC 2 Audit Trail Completeness}
\index{SOC 2}
\index{audit trail}

For SOC 2 Type II assessments, auditors typically request evidence of:

\begin{description}
  \item[Access control decisions] Every connection produces a JSON log entry with
    action, reason, and client identification. The structured log format provides
    a complete, queryable audit trail.
  \item[Configuration change management] The \code{management:policy\_audit} Redis
    list records every policy change with timestamp, operator attribution, and
    old/new values. This list is limited to 1000 entries — for long-term retention,
    ship audit log entries to an external store.
  \item[Security monitoring] Prometheus metrics, Grafana dashboards, and Alertmanager
    alert rules provide evidence of continuous security monitoring.
  \item[Incident response] The analytics node's campaign detection and score drift
    alerting provide evidence of automated anomaly detection.
\end{description}

\section{GDPR Data Flow Diagram}
\index{GDPR!data flow}

\begin{figure}[h]
\begin{fullwidth}
\begin{tikzpicture}[
  node distance=1.2cm,
  box/.style={draw=ja4navy, fill=ja4lightgray, rounded corners=3pt,
              minimum width=3cm, minimum height=0.8cm,
              font=\footnotesize\sffamily, text centered},
  store/.style={draw=ja4purple, fill=ja4purplelight, rounded corners=3pt,
                minimum width=3cm, minimum height=0.8cm,
                font=\footnotesize\sffamily, text centered},
  ext/.style={draw=ja4gray, fill=white, rounded corners=3pt,
              minimum width=2.8cm, minimum height=0.8cm,
              font=\footnotesize\sffamily, text centered},
  arrow/.style={->, >=Stealth, draw=ja4navy},
  label/.style={font=\tiny\sffamily\itshape, text=ja4gray},
]
  \node[box] (client) {Data Subject\\(Client IP)};
  \node[box, right=1.5cm of client] (proxy) {\ja\ Proxy\\(Controller's system)};
  \node[store, right=1.5cm of proxy] (redis) {Redis\\(IP + metadata\\TTL varies)};
  \node[store, above=0.8cm of redis] (loki) {Loki\\(Logs with IP\\30-day retention)};
  \node[store, below=0.8cm of redis] (prom) {Prometheus\\(Metrics\\15-day default)};
  \node[ext, right=1.5cm of redis] (abuseipdb) {AbuseIPDB\\(3rd party processor)};
  \node[ext, above=0.8cm of abuseipdb] (rdap) {RDAP / RIR\\(Public data)};
  \node[ext, below=0.8cm of abuseipdb] (maxmind) {MaxMind\\(GeoLite2\\local DB)};

  \draw[arrow] (client) -- node[above, label]{IP addr} (proxy);
  \draw[arrow] (proxy) -- node[above, label]{ban/visit/sess} (redis);
  \draw[arrow] (proxy) -- node[left, label]{JSON log} (loki);
  \draw[arrow] (proxy) -- node[left, label]{metrics} (prom);
  \draw[arrow] (proxy) -- node[above, label]{IP lookup} (abuseipdb);
  \draw[arrow] (proxy) -- node[above, label]{IP lookup} (rdap);
  \draw[arrow] (maxmind) -- node[below, label]{country/ASN} (proxy);
\end{tikzpicture}
\end{fullwidth}
\caption{GDPR data flow showing all systems where client IP addresses are processed or stored.}
\label{fig:gdpr-flow}
\end{figure}

\subsection{Third-Party Data Processors}

\ja\ sends client IP addresses to the following third-party services. Operators must
ensure appropriate data processing agreements (DPAs) are in place:

\begin{fullwidth}
\begin{table}[h]
\caption{Third-party data processors}
\label{tab:processors}
\begin{tabular}{@{}lllp{4cm}@{}}
\toprule
\textbf{Service} & \textbf{Data sent} & \textbf{Retention by processor} & \textbf{DPA required?} \\
\midrule
AbuseIPDB & Client IP address & Per AbuseIPDB privacy policy & Yes, if processing EU personal data \\
RDAP/WHOIS (RIRs) & Client IP address & Public records; varies by RIR & Assess per RIR \\
MaxMind GeoLite2 & None (local DB) & N/A — local file, no network calls & No \\
DNS resolver & Reverse DNS query for client IP & Per resolver retention policy & Depends on resolver \\
\bottomrule
\end{tabular}
\end{table}
\end{fullwidth}

\begin{notebox}[Disabling AbuseIPDB for Privacy Compliance]
If your jurisdiction or DPA requires that client IP addresses not be sent to third
parties without consent, set \configkey{abuseipdb.enabled: false}. This removes the
only signal that requires active transmission of client IPs to an external service.
RDAP queries can also be disabled (\configkey{rdap.enabled: false}).
MaxMind GeoLite2 operates entirely on a local database file with no network calls.
\end{notebox}
